\section{Background}

\subsection{Intelligent Textbooks and MicroSims}

An intelligent textbook differs from a traditional digital textbook in several
structural ways. It is typically organized around a \textbf{concept dependency
graph}---a directed acyclic graph in which nodes represent individual concepts
and edges represent prerequisite relationships. Learning objectives are tagged
against a taxonomy such as Bloom's Taxonomy~\cite{bloom1956taxonomy}, which
orders cognitive demand from basic recall (\textit{Remember},
\textit{Understand}) through application and analysis (\textit{Apply},
\textit{Analyze}) to synthesis and judgment (\textit{Evaluate}, \textit{Create}).

The most distinctive feature is the \textbf{MicroSim}: a small, self-contained
interactive simulation that lets a learner manipulate parameters and observe
outcomes~\cite{lockhart2025microsims}. A MicroSim built with the p5.js
JavaScript library, for example, might let a student adjust the initial
velocity of a bouncing ball and observe the resulting trajectory, building
intuition for Newtonian mechanics. MicroSims are typically packaged as a
directory containing a main HTML file, the simulation JavaScript, and
supporting metadata, and are embedded into textbook pages via an inline frame.

A useful design principle for intelligent textbooks is the
\textit{five-level framework}, developed in detail by
McCreary~\cite{mccreary2025fivelevels}. Table~\ref{tab:five-levels} presents
the five levels, using the names established in that framework, together with
two threshold columns that carry significant practical implications for
institutional deployment.

\begin{table}[h]
\centering
\small
\caption{Five Levels of Intelligent Textbook Capability
  (names after~\cite{mccreary2025fivelevels})}
\begin{tabular}{@{}cL{2.1cm}L{4.7cm}C{1.4cm}C{1.4cm}@{}}
\toprule
Level & Name & Key Institutional Requirements & Student Data & LLM Required \\
\midrule
1 & Static      & Fixed content only; no institutional systems required          & No  & No  \\
2 & Interactive & Static hosting; anonymous analytics optional                   & No  & No  \\
3 & Adaptive    & Student authentication; institutional data infrastructure;
                  ongoing relationship management                                & Yes & No  \\
4 & Chatbot     & Integration with institutional systems; LLM oversight
                  capabilities; trust relationships                              & Yes & Yes \\
5 & Autonomous  & Deep institutional knowledge; comprehensive student
                  support systems                                               & Yes & Yes \\
\bottomrule
\end{tabular}
\label{tab:five-levels}
\end{table}

Two threshold boundaries emerge from Table~\ref{tab:five-levels} with
significant practical implications. The first is a \textbf{data-privacy
boundary} between Levels~2 and~3: Levels~1 and~2 require no student-specific
data and can be served to completely anonymous users, while the Adaptive
level (Level~3) and above demand student authentication, individual learning
histories, assessment records, and ongoing institutional relationship
management. This boundary triggers obligations under data-protection
regulations such as FERPA, COPPA, and GDPR, and it is the threshold below
which institutions can deploy intelligent textbooks with minimal governance
overhead.

The second is an \textbf{LLM-cost boundary} between Levels~3 and~4.
Levels~1 through~3 can be implemented entirely with deterministic
algorithms---graph traversal, rule-based sequencing, and spaced
repetition---without any large language model inference. A well-constructed
concept dependency graph, combined with per-learner mastery estimates, is
sufficient to drive adaptive navigation without generative AI. The Chatbot
level (Level~4) and the Autonomous level (Level~5), by contrast, require
continuous LLM access---Level~4 for real-time conversational tutoring and
Level~5 for autonomous content generation and self-modification. This
dependency increases operating costs by one to two orders of magnitude
relative to lower levels and introduces ongoing dependencies on third-party
model providers.

This paper is fundamentally concerned with the transition from the Chatbot
level (Level~4) to the Autonomous level (Level~5)---the point at which the
textbook begins to improve its own content based on empirical evidence from
learning analytics. Because Level~5 inherits the LLM dependency of Level~4
and adds the cost of continuous fitness evaluation and variant generation,
the evolutionary loop proposed here must account for inference cost as a
component of overall system viability.

\subsection{The Darwin-G\"{o}del-Machine}

The Darwin-G\"{o}del-Machine is a self-improving system that iteratively
modifies its own code while empirically validating each change against coding
benchmarks~\cite{zhang2025dgm}. Its design draws on two intellectual
traditions. From the theoretical G\"{o}del Machine it inherits the goal of
self-referential improvement, but it replaces the requirement of formal proof
with empirical benchmark validation. From Darwinian evolution and
open-endedness research it inherits the practice of maintaining an
\textit{archive} of generated solutions rather than evolving a single line of
descent.

The system operates in an iterative cycle. It begins with a base coding agent
built on a frozen foundation model equipped with tool-use capabilities. In
each generation, a parent agent is selected from the archive, the parent is
tasked with modifying itself to produce a child variant, the child is
evaluated on benchmarks, and successful children are added to the archive. The
archive enables open-ended exploration: previously suboptimal designs can serve
as stepping stones to later breakthroughs, and many distinct improvement paths
can be explored in parallel.

Reported results are substantial. The DGM improved its performance on the
SWE-bench coding benchmark from 20.0\% to 50.0\%, and on the Polyglot
multi-language benchmark from 14.2\% to 30.7\%. Ablation studies demonstrated
that both self-modification and open-ended exploration were essential;
disabling either substantially degraded performance. The authors emphasize that
all experiments were conducted with safety precautions including sandboxing and
human oversight.

Three properties of the DGM are directly relevant to educational content:

\begin{enumerate}
    \item \textbf{Empirical over formal validation.} We cannot prove a MicroSim
          is pedagogically optimal, but we can measure whether students learn
          better from one variant than another.
    \item \textbf{Archival open-endedness.} Maintaining a population of content
          variants avoids premature convergence on a locally good but globally
          suboptimal design.
    \item \textbf{Frozen foundation, evolving scaffold.} The DGM does not
          modify its underlying model; it modifies the code and tools around it.
          Analogously, an educational system need not modify the underlying
          language model---only the content artifacts and the prompts that
          generate them.
\end{enumerate}

\subsection{The Broader Landscape of Self-Improving Systems}

The DGM is one of several recent systems exploring self-improvement, and
understanding the landscape clarifies which mechanisms are most applicable to
education.

\textbf{AlphaEvolve}~\cite{novikov2025alphaevolve} (Google DeepMind, May 2025)
is an evolutionary coding agent that pairs large language models with
evolutionary computation to discover and refine algorithms. It orchestrates an
autonomous pipeline of LLMs that propose direct changes to code, with one or
more automatic evaluators providing feedback that grounds the evolutionary
process. AlphaEvolve requires an evaluation function with metrics to optimize
and an initial algorithm to seed the search. Among its notable results, it
discovered a procedure to multiply two $4 \times 4$ complex-valued matrices
using 48 scalar multiplications, the first improvement over Strassen's 1969
algorithm in that setting, and it produced practical optimizations to data
center scheduling and hardware accelerator circuit design. AlphaEvolve's
reliance on an explicit, automatic evaluation function is the single most
transferable idea for educational content: the evaluator \textit{is} the
fitness function, and designing it well is the central challenge.

\textbf{SEAL (Self-Adapting Language Models)}~\cite{zweiger2025seal} (MIT,
June 2025; expanded for NeurIPS 2025) takes a different approach, enabling a
model to generate its own fine-tuning data and update directives. The model
produces ``self-edits''---natural-language descriptions of what should
change---which are converted into fine-tuning examples and applied via
supervised learning, with a reinforcement learning outer loop that uses
downstream task performance as the reward signal. On a factual
question-answering task, SEAL improved accuracy from 32.7\% to 47.0\% after
two rounds of training, outperforming fine-tuning on raw passages or on
synthetic data generated by a strong external model. SEAL's relevance to
education is its focus on \textit{knowledge incorporation}---the same problem a
textbook faces when it must integrate new material---and its explicit
acknowledgment that human supervision remains necessary to ensure improvements
are beneficial and aligned.

\textbf{Open-source successors} have followed quickly.
CodeEvolve~\cite{assuncao2025codeevolve} reproduces the high-level principles
of AlphaEvolve in an open framework using an island-based genetic algorithm
with a weighted LLM ensemble, reporting competitive performance at a fraction
of the compute cost. The G\"{o}del Agent~\cite{yin2024godelagent} sketches a
self-referential architecture in which an agent proposes and accepts
modifications to itself based on environmental feedback.

Table~\ref{tab:systems} summarizes the landscape.

\begin{table}[h]
\centering
\small
\caption{Landscape of Self-Improving AI Systems and Applicability to Textbooks}
\begin{tabular}{@{}L{2.6cm}L{2.4cm}L{2.0cm}L{2.4cm}L{2.4cm}@{}}
\toprule
System & Organization & What evolves & Validation signal & Textbook applicability \\
\midrule
G\"{o}del Machine & Schmidhuber (theory) & Own code & Formal proof & Low (intractable) \\
Darwin-G\"{o}del-Machine & Sakana AI / UBC / Vector & Agent code + tools & Coding benchmarks & High (archival model) \\
AlphaEvolve & Google DeepMind & Algorithms / code & Automatic evaluators & High (fitness function) \\
SEAL & MIT & Model weights & Downstream task RL reward & Medium (knowledge incorporation) \\
CodeEvolve & Open source & Code & Benchmark evaluators & High (reproducible, low cost) \\
\bottomrule
\end{tabular}
\label{tab:systems}
\end{table}

The consistent pattern across all successful systems is the combination of a
generative engine (a foundation model proposing variants) with a grounding
evaluator (an automatic measure of quality). For educational content, the
generative engine already exists in the form of code-generation tools; the
missing piece is a well-designed evaluator built from learning analytics. The
rest of this paper focuses on building that evaluator.
